\chapter{预训练：构建可复用的数据表示}
\label{chap:pretraining}
\begin{quote}\small
\textit{本章摘要。}预训练通过先行学习任务获得可复用的参数或表示。本章从数据与目标的关系出发，解释监督预训练、掩码与生成式学习、对比学习、无负样本自蒸馏及领域继续预训练。工业振动信号案例用于说明如何选择增强与遮蔽方式，以及如何区分预训练损失下降、表示改善和下游收益。
\end{quote}

前面的模型章节回答序列、图像和文本应如何表示，上一章进一步让物理方程参与学习：当标签不足时，已知规律可以约束候选解。然而，不是所有设备都有足够准确的控制方程，许多可利用的知识只隐含在历史观测中。本章由此转向另一条路线：在目标故障标签不足时，能否先利用历史数据学习可复用结构，再用有限标签建立预测器？预训练的价值在于可复用性，其监督信号可以来自人工标签，也可以来自数据本身。它不等同于无监督学习，更不要求模型一定具有很大规模。

理解这条路线，需要区分三个并不等价的问题：辅助任务能否被完成，表示是否保留了任务需要的信息，有限标签能否把这些信息转成可靠判断。重构波形可以依靠平滑性，对比学习可以依靠设备身份，而故障诊断可能恰恰需要低能量的短时冲击。因此，本章不以罗列预训练模型为主线，而是逐步追问学习信号要求模型保留什么、丢弃什么：先由似然与重构解释预测目标，再由正负样本关系解释不变性，最后分析不使用负样本时如何避免所有表示趋同。贯穿这些方法的观点是，预训练目标是一种关于有用结构的假设，必须由独立的下游任务检验。

读完本章，可以按三个问题检查自己的方案：哪些历史记录允许参与训练，遮住的部分是否太容易猜中，模型是否把不同输入都编码成了相近的结果。例如波形补得很平滑，却丢掉短时冲击，就未必有利于故障诊断。下一章继续讨论拿到预训练参数以后，哪些可以保留，哪些需要用新设备的数据重新调整。

\section{预训练学到了什么}

% auxiliary-resource-guide
本章末的“辅助学习材料”列出与核心方法对应的讲解和实践入口，可在阅读相关推导后，按所列小节对照学习。

\begin{figure}[htbp]
\centering
\begin{tikzpicture}[node distance=0.75cm and 0.75cm]
 \node[idi input, minimum width=2.0cm] (data) {大规模数据\\$\mathcal D_{\rm pre}$};
 \node[idi process, right=of data, minimum width=2.0cm] (obj) {预训练目标\\重构/对比/生成};
 \node[idi state, right=of obj, minimum width=2.0cm] (enc) {可复用编码器\\$f_{\theta_0}$};
 \node[idi output, right=of enc, minimum width=2.0cm] (task) {目标任务\\线性探测或微调};
 \draw[idi arrow] (data)--(obj); \draw[idi arrow] (obj)--(enc); \draw[idi arrow] (enc)--(task);
 \draw[idi feedback] (task.south) .. controls +(0,-0.42) and +(0,-0.42) .. node[below,font=\scriptsize]{下游收益反查预训练目标} (obj.south);
\end{tikzpicture}
\caption{预训练的价值由下游任务检验：辅助目标下降并不自动意味着故障相关信息被保留。}
\label{fig:pretraining-downstream-validation}
\end{figure}
图\ref{fig:pretraining-downstream-validation}把无标签预训练与下游评估连成反馈路径，强调辅助任务学到的表示须通过独立故障任务检验，而不能只用预训练损失验收。反馈只允许使用训练与验证任务：若根据测试集表现反复改变预训练目标，该测试集就参与了模型选择，不能再充当最终独立证据。
阅读下面的目标函数，只需先掌握小批次训练、交叉熵和反向传播。与从零训练相比，本章多了一个先行任务：用什么训练信号产生可复用参数？先辨认每个目标的输入和答案来自哪里，再判断学到的特征能否帮助有限标签任务。

\paragraph{表示、编码器和头分别是什么。}
表示是从原始输入计算出的特征向量，而不是人工故障标签。例如一个窗口含$C=3$个通道、每通道$L=100$个采样值，可展平为$x\in\R^{300}$，编码器$f_\theta:\R^{300}\to\R^{32}$把它变成$32$维表示。重构头$g_\phi:\R^{32}\to\R^{300}$把表示映射回波形；下游三分类头则把同一表示映射到$3$个类别分数，经 softmax 变成概率。头是为特定目标服务的末端模块，骨干或编码器负责提取共享特征。不同任务的头输出尺寸不同，所以预训练结束时常需替换它，而非直接把重构波形当成故障概率。

在线性探测中，固定$h=f_{\theta_0}(x)\in\R^{32}$，只学习$W\in\R^{3\times32}$、$b\in\R^3$，计算类别分数$a=Wh+b$和概率$\softmax(a)$。这里“线性”指分数对表示$h$是仿射的，不是对原始波形一定线性，也不是把 softmax 删除。若表示已丢掉故障冲击，再训练这样的头不能恢复被删信息；全量微调则允许编码器重新学习，但需要更多目标证据控制过拟合。

“预训练”描述先学习可复用参数的阶段，“自监督”描述标签如何由数据本身构造，“迁移学习”描述把已有知识用于目标问题的做法。它们不互斥：用掩码恢复历史波形是自监督预训练，把得到的编码器用于新设备故障分类是迁移，训练新分类头或更新编码器则是具体适配操作。监督预训练也能用于迁移，迁移也不必使用单独设计的自监督阶段。

设预训练数据为$\mathcal D_{\mathrm{pre}}$，编码器为$f_\theta$，与预训练目标配套的预测头为$g_\phi$。一般训练形式为
\begin{equation}
 (\theta_0,\phi_0)=\argmin_{\theta,\phi}
 \E_{x\sim\mathcal D_{\mathrm{pre}}}
 \E_{\omega\sim P_\omega}
 \ell_{\mathrm{pre}}\bigl(g_\phi(f_\theta(v_\omega(x))),s_\omega(x)\bigr).
\end{equation}
式中$\argmin$表示寻找使平均损失尽量小的参数，外层期望平均不同记录，内层期望平均同一记录的随机处理方式，实际通常以一个小批次近似。下标$0$表示交付下游任务的预训练参数，并非所有参数取零。监督版本的外层抽样应读作$(x,y)\sim\mathcal D_{\mathrm{pre}}$，此时目标$s_\omega$取该记录的人工标签$y$；不能只凭输入$x$凭空得到人工标签。

其中$\omega$控制裁剪、遮蔽或目标构造，$v_\omega(x)$为模型输入，$s_\omega(x)$为学习目标。监督分类时目标来自人工标签，自监督任务则由样本结构构造目标，例如被遮蔽的片段或同一样本的另一视图。预训练完成后可保留编码器、丢弃辅助头，再接入目标任务头$h_\psi$。

若冻结$\theta_0$只训练线性$h_\psi$，检验的是该表示中的线性可分信息；全量微调则同时检验初始化能否支持适应。两者可能产生不同排序。模型也可能保留了设备编号和工况，却没有保留故障机理，因此预训练损失不能直接用作表示质量指标。下一章讨论微调、参数高效适配与域对齐，本章重点是初始化产生之前的学习过程。

预训练数据不一定与下游任务同分布。通用图像表示可能提供边缘和纹理结构，但工业细裂纹要求保留更细的空间信息；自然语言模型能够学习句法，却未必知道某台设备的有效规程。数据规模、任务相关性和目标设计共同决定迁移潜力，增加其中一项不能保证补偿另一项的缺失。

\section{监督预训练与重建式目标}
监督预训练利用较大源数据集的标签训练骨干，例如图像分类网络中的交叉熵目标\cite{he2016resnet}：
\begin{equation}
 \mathcal L_{\mathrm{sup}}=-\frac1N\sum_{i=1}^{N}
 \log p_{\theta,\phi}(y_i\mid x_i).
\end{equation}
随后保留骨干并替换分类头。源标签可以来自其他设备或相关任务，但其类别定义决定了哪些差异受到强调。只以设备类别预训练，可能更容易区分传感器安装位置而不是早期故障。弱标签和伪标签能扩大覆盖，却需要单独处理标签噪声，不能把它们等同于准确人工标注。

自编码器从压缩表示重构输入，去噪自编码器则从受扰动输入$\widetilde x$恢复原始$x$\cite{vincent2008denoising}：
\begin{equation}
 \mathcal L_{\mathrm{denoise}}
 =\E_{x\sim\mathcal D_{\mathrm{pre}}}\E_{\widetilde x\sim q(\cdot\mid x)}
 \|g_\phi(f_\theta(\widetilde x))-x\|_2^2.
\end{equation}
扰动机制决定模型需要恢复什么。噪声过小可能只要求近似恒等映射，噪声过大可能使恢复目标含糊。重构均方误差倾向于优先解释高能量成分，低幅值但诊断价值较高的冲击不一定受到足够重视。频谱或多尺度损失可以改变关注点，但要明确引入的频段先验，避免利用测试故障选择损失。

\subsection{从条件似然理解重构为何产生平滑结果}
把一个窗口展平为$x\in\R^D$，令$u=\widetilde x$表示受扰动观测，编码器输出$f_\theta(u)\in\R^d$，解码器输出$m_{\theta,\phi}(u)\in\R^D$。若建模假设为
\begin{equation}
 p_{\theta,\phi}(x\mid u)=\mathcal N(x;m_{\theta,\phi}(u),\sigma^2 I_D),
\end{equation}
则单个样本的负对数似然是
\begin{equation}
 -\log p(x\mid u)=\frac D2\log(2\pi\sigma^2)
 +\frac1{2\sigma^2}\|x-m_{\theta,\phi}(u)\|_2^2.
\end{equation}
只有在$\sigma^2$固定时，最大似然才等价于普通均方误差最小化；若同时学习方差，就不能丢弃对数方差项。不同通道单位差异很大时，采用对角协方差$\Sigma$会得到$(x-m)^\top\Sigma^{-1}(x-m)$，即按噪声尺度加权的误差。这样的权重表达的是可靠性或尺度假设，不能任意解释为故障重要性。

进一步令$m^*(u)=\E[X\mid U=u]$，对任意候选重构$a$展开平方：
\begin{align}
 \E[\|X-a\|^2\mid u]
 &=\E[\|X-m^*+m^*-a\|^2\mid u]\\
 &=\E[\|X-m^*\|^2\mid u]+\|m^*-a\|^2.\label{eq:pretraining-conditional-mse}
\end{align}
式\eqref{eq:pretraining-conditional-mse}中交叉项消失，是因为$\E[X-m^*\mid u]=0$。因此无限容量、总体风险最优的平方误差预测器输出条件均值。假如遮蔽区的冲击可能出现在两个不同位置，条件均值可能在两处各产生半个冲击，虽然平均误差最小，却不是任何一次真实波形。这不是单纯增加网络容量就能解决的问题；需要缩短遮蔽区、加入可靠上下文，或用能表达多峰条件分布的预测头。

用两个等可能的真实片段$(2,0)$与$(0,2)$可检查条件均值结论。预测$(1,1)$时，两种情况的平方误差都为$2$；固定预测$(2,0)$时，误差分别为$0$和$8$，平均为$4$。因此均值确实赢得平方损失，却把一个完整冲击拆成两个半冲击。重构指标和故障语义之间的差异由此可见。

\subsection{遮蔽任务与因果预测的边界}
掩码建模只显示部分输入，并预测被遮蔽区域。令$M$为目标位置集合、$V$为可见位置集合，可写为
\begin{equation}
 \mathcal L_{\mathrm{mask}}
 =\E_{x,M}\frac1{|M|}\sum_{j\in M}
 \ell\bigl(\widehat x_j(x_V,M),x_j\bigr).
\end{equation}
离散词元通常采用交叉熵，连续图像块或信号片段可采用重构误差。BERT 使用掩码语言建模，MAE 则以可见图像块训练编码器，并由轻量解码器重构被遮蔽块\cite{devlin2019bert,he2022mae}。MAE 的编码器只处理可见块，减少了高遮蔽率下的编码开销；解码器仍需接收位置信息与遮蔽标记。不能仅按遮蔽比例推断所有架构都具有相同节省。

这里假定输入由$L$个元素或片段组成，每个连续片段$x_j\in\R^p$，$M\subset\{1,\ldots,L\}$非空，$V$为其补集；解码器需为每个$j\in M$输出同维预测。若采用逐位置条件分布$q_{\theta,\phi}(x_j\mid x_V,M)$，掩码损失就是条件负对数似然之和。逐位置相加相当于用条件独立输出头拟合各位置的边缘分布，不表示真实被遮蔽片段给定可见输入后必然独立；若要生成相互协调的一整段，还需联合或自回归解码器。$1/|M|$避免遮蔽更多位置时仅因项数增加而提高损失，但它不会消除不同遮蔽难度带来的梯度差异。

离散监督分类和掩码词元分类共享交叉熵结构。令真实条件分布为$r(k\mid u)$、预测为$p(k\mid u)$，则$-\sum_kr_k\log p_k=H(r)+\KL(r\|p)$。固定输入条件下，总体最优预测应匹配$r$；有限标签训练以单个类别指示向量代替未知$r$。因此监督预训练强调源标签可区分的信息，掩码训练强调由上下文可预测的信息，二者都没有数学理由自动保留与自身目标无关的下游细节。

例如$x=(1,2,3,4)$，遮蔽集合$M=\{2,3\}$，输入只显示位置$1,4$及其数值。若模型对遮蔽位置输出$(2.5,2.5)$，则按本节平方误差约定，损失为$[(2.5-2)^2+(2.5-3)^2]/2=0.25$；可见位置不参与这次目标平均。监督值来自原始完整记录，但不能先把完整记录编码后才遮蔽目标，否则模型已经看过答案。连续片段含$p$个数时，单位置损失可用平方范数，是否再除以$p$也应统一。

独立随机点遮蔽可能被局部平滑插值轻易解决。时间序列可使用连续片段遮蔽、跨通道遮蔽或多尺度遮蔽，以迫使模型利用更长上下文。遮蔽过长又可能使目标不可预测，网络只能输出条件平均。应根据信号周期、传感器耦合与任务时间尺度选择策略，并检查模型是否只是复制相邻样本。

自回归预训练预测给定过去的下一个元素，目标为
\begin{equation}
 \mathcal L_{\mathrm{AR}}=-\E_x\sum_{j=1}^{L}
 \log p_\theta(x_j\mid x_{<j}).
\end{equation}
该目标来自联合分布的链式分解$p_\theta(x_1,\ldots,x_L)=\prod_{j=1}^Lp_\theta(x_j\mid x_{<j})$；取负对数将连乘转为求和，不要求序列元素相互独立。训练时使用真实前缀计算条件似然，而生成时使用模型已生成的前缀，两者输入分布可能不同。对总体数据分布$P$，期望负对数似然等于$H(P)+\KL(P\|P_\theta)$；模型族足够表达且优化达到总体最优时，才有匹配数据分布的解释。有限数据、有限上下文或模型失配下，即便整体似然改善，稀有故障的条件概率也可能仍很差。对文本这是下一词元预测，对连续信号需要明确预测分布或回归损失。双向掩码表示可以利用目标位置两侧上下文，自回归表示受因果边界约束。在线故障预测若采用窗口表示，应保证窗口中没有预测时刻之后的数据，不能因任务名称为自监督就忽略时间泄漏。

\section{对比学习：哪些样本应当接近}
对比学习用正样本对定义应保留的一致性，用负样本提供区分信号。A Simple Framework for Contrastive Learning of Visual Representations (SimCLR) 对一个批次中的每个样本独立增强两次，经编码器与投影头得到归一化表示$z_i$\cite{chen2020simclr}。对锚点$i$及其正样本$j$，损失为
\begin{equation}
 \ell_{i,j}=-\log
 \frac{\exp(z_i^\top z_j/\tau)}
 {\sum_{k\ne i}\exp(z_i^\top z_k/\tau)},\qquad \tau>0.
\end{equation}
含$B>1$个原始样本的批次产生$2B$个视图；若$B=1$，只有正候选，损失恒为零，无法提供这里的对比信号。分母含一个正样本和$2B-2$个负样本，训练对两个方向求平均。投影头使对比目标作用于专门的空间，下游任务常使用投影前的编码器表示。温度$\tau$改变相似度差异对梯度的影响，较小温度强调难区分样本，也可能放大假负样本的影响。

“视图”就是同一条记录经不同允许处理得到的输入，正样本不是“正常设备”，负样本也不是“故障设备”，而是当前配对任务中匹配与不匹配的候选。例如$B=2$，某锚点有一个正候选和两个负候选，若与三者的相似度分别为$1,0,0$且$\tau=1$，则正候选概率为$e/(e+2)\approx0.576$，损失约$0.551$。若三者都同样相似，则概率为$1/3$、损失为$\log3\approx1.099$。这说明对比训练不是单独提高正相似度，而是提高它相对于其他候选的辨识度。

\paragraph{把对比损失展开为一个分类问题。}
设投影前表示为$h_i\in\R^d$，投影后归一化向量为$z_i\in\R^m$，$\|z_i\|_2=1$。固定锚点$i$时，其余$2B-1$个视图构成候选集合。令$s_{ik}=z_i^\top z_k$及
\begin{equation}
 q_{ik}=\frac{e^{s_{ik}/\tau}}{\sum_{l\ne i}e^{s_{il}/\tau}}.
\end{equation}
任务就是从候选中识别真正配对的视图$j$。把$\ell_{i,j}=-s_{ij}/\tau+\log\sum_{k\ne i}e^{s_{ik}/\tau}$逐项求导，有
\begin{equation}
 \frac{\partial\ell_{i,j}}{\partial s_{ik}}
 =\frac1\tau(q_{ik}-\mathbf1[k=j]),\qquad
 \nabla_{z_i}\ell_{i,j}
 =\frac1\tau\left(\sum_{k\ne i}q_{ik}z_k-z_j\right).
\end{equation}
这里第二式暂把其他向量视为常量；完整批次梯度还包含$z_i$作为别的锚点的候选时收到的梯度。下降方向拉近正样本并远离概率加权的候选均值，而非均匀排斥每个负样本。若未归一化投影为$u_i\ne0$，还需乘归一化雅可比$(I-z_iz_i^\top)/\|u_i\|_2$，不能把单位球面上的变量当成任意自由向量。零向量不能直接归一化；实现中若用小常数保护分母，保护区间内的导数也应按实际公式计算。

这个推导解释了温度的两个作用：显式的$1/\tau$改变梯度尺度，softmax 又把权重集中到最相似候选。因此减小温度不只是“加大学习率”。当一个高相似负样本实际属于同类故障时，它会得到很大的排斥权重。另一方面，若所有$z$完全相等，则$q_{ik}=1/(2B-1)$，损失为$\log(2B-1)$，没有完成配对识别；但对称构型仍可能有退化驻点，不能仅从负样本存在推断优化必然成功。

大批次可以提供更多负样本，但增加内存成本。Momentum Contrast (MoCo) 使用键队列复用近期表示，并用动量编码器降低队列中新旧表示的不一致\cite{he2020moco}。若查询编码器参数为$\theta_q$、键编码器参数为$\theta_k$，则
\begin{equation}
 \theta_k\leftarrow\mu\theta_k+(1-\mu)\theta_q,\qquad 0\leq\mu<1.
\end{equation}
查询分支通过梯度更新，键分支按指数移动平均更新，当前键入队并移除最旧键。队列增大并不总有益：分布变化较快时，陈旧键可能形成不合适的比较对象。

正样本定义比损失名称更重要。图像颜色扰动可以弱化照明差异，但在变色故障中颜色本身就是标签信息。振动信号的强幅值归一化可能删除故障严重度，时间伸缩可能改变与转速相关的频率结构。只有在目标任务对某种变化确实不敏感时，才应训练这种不变性。状态变化对应的敏感方向应保留，而不是一律压缩。

同一故障的两个不同记录可能被当作负样本，造成假负样本排斥。监督对比学习利用已有类别标签为一个锚点构造多个正样本\cite{khosla2020supcon}，但其适用范围由标签覆盖决定。工业采样还应控制设备和批次：若一个批次中每台设备只出现一种工况，模型可能靠设备身份完成对比任务，而非学习可迁移的状态表示。

跨模态对比预训练将同步波形、图像或文本对齐。CLIP 使用图像与文本的双向批次对比目标\cite{radford2021clip}，其思想可以扩展到工业配对数据，但工单时间不一定等于故障发生时间，粗糙配对会引入错误正样本。跨模态目标学到的是给定配对规则下的关联，不会自动建立因果对应。

\section{无负样本学习与表示坍缩}
如果只令两个增强视图的表示接近，所有输入都映射到同一个常量也能使距离为零，这称为表示坍缩。Bootstrap Your Own Latent (BYOL) 通过在线编码器、预测头和移动平均目标编码器构成不对称训练\cite{grill2020byol}。用$\operatorname{sg}$表示停止梯度，归一化后的预测与目标分别为$\overline q_\theta$和$\overline z_\xi$，则一个方向的损失为
\begin{equation}
 \ell_{\mathrm{BYOL}}=
 \left\|\overline q_\theta(v_1)
 -\operatorname{sg}\bigl(\overline z_\xi(v_2)\bigr)\right\|_2^2,
 \qquad \xi\leftarrow\mu\xi+(1-\mu)\theta.
\end{equation}
实际训练交换视图后再平均。这里$\theta$和$\xi$表示两分支中对应编码与投影模块的参数，在线预测头单独通过梯度更新。停止梯度和移动平均控制目标变化，并不使任意网络与任意数据增强都自动免于坍缩。

Simple Siamese (SimSiam) 去掉移动平均目标更新，使用共享编码器、预测头和停止梯度构造更简洁的孪生学习\cite{chen2021simsiam}。Self-Distillation with No Labels (DINO) 则让学生匹配教师在不同裁剪视图上的概率输出，并以移动平均教师、输出中心化和温度锐化调节目标分布\cite{caron2021dino}。这些方法的教师通常来自学生训练轨迹，而不是额外人工标签，因此属于自蒸馏。训练中要检查特征方差、有效秩和样本间距离，不能只看两个分支越来越一致。

另一类方法直接约束表示统计量。VICReg 组合视图一致性、各维方差下限和非对角协方差惩罚\cite{bardes2022vicreg}：
\begin{equation}
 \mathcal L=\lambda_s\mathcal L_{\mathrm{inv}}
 +\lambda_v\mathcal L_{\mathrm{var}}
 +\lambda_c\mathcal L_{\mathrm{cov}}.
\end{equation}
对中心化批次表示$Z\in\R^{B\times d}$，协方差为$C=Z^\top Z/(B-1)$，协方差项惩罚$C$的非对角元素平方，方差项则惩罚标准差低于给定阈值的维度。两项通常分别作用于两个视图分支。它以显式统计约束抑制常量表示和维度冗余，但小批次统计不稳定，低相关性也不保证所有维度都具有诊断意义。

\paragraph{停止梯度和统计约束究竟改变了什么。}
对单位向量$q,z$，$\|q-z\|^2=2-2q^\top z$，所以 BYOL 的距离最小化也在最大化方向一致性。若目标$z$停止梯度，单次更新只改变在线分支，而目标编码器通过移动平均间接改变。把动量递推展开，得到
\begin{equation}
 \xi_t=\mu^t\xi_0+(1-\mu)\sum_{s=1}^t\mu^{t-s}\theta_s.
\end{equation}
教师因此是历史学生的指数加权组合：较大的$\mu$平滑噪声，也增加对新分布的响应延迟。停止梯度不是修改距离的数值，而是修改计算图的导数路径；它与预测头、归一化和优化共同作用，单独的常量映射依然能让一致性项为零。

为了明确 VICReg 如何额外排斥这个解，令未中心化批次表示为$Z^{(a)}\in\R^{B\times d}$、$a\in\{1,2\}$，$B>1$，中心化后协方差为$C^{(a)}$。一组明确的损失约定是
\begin{align}
 \mathcal L_{\mathrm{inv}}&=\frac1B\sum_{i=1}^B\|z_i^{(1)}-z_i^{(2)}\|_2^2,\\
 \mathcal L_{\mathrm{var}}&=\frac1{2d}\sum_{a=1}^2\sum_{j=1}^d
 \max\{0,\gamma-\sqrt{C_{jj}^{(a)}+\epsilon}\},\\
 \mathcal L_{\mathrm{cov}}&=\frac1{2d}\sum_{a=1}^2\sum_{j\ne k}(C_{jk}^{(a)})^2.
\end{align}
这里$\gamma>\sqrt\epsilon>0$规定标准差下限。常量表示使一致性和协方差项都为零，却令每一维方差惩罚严格为正；所有维度复制同一个变化信号能够通过方差检查，却会产生很大的非对角协方差。这说明两种统计惩罚分别针对无变化和重复编码，缺一不可。不过$\operatorname{rank}(C)\leq\min(d,B-1)$，当$d>B-1$时不可能在一个批次内同时达到全维正方差与严格零非对角协方差。损失权重是在这些目标间折中，而不是保证有限批次出现理想白化表示。

可把停止梯度先看成标量操作：$L=(a-\operatorname{sg}(b))^2$在$a=1,b=2$时值为$1$，对$a$的导数为$-2$，而沿被停止的目标路径对$b$的导数为零。若移去$\operatorname{sg}$，数值仍为$1$，但对$b$的导数变为$2$。教师不是“知道真相”的专家，只是提供当前更新暂不追赶的目标。比如教师参数$\xi=0$、学生对应参数$\theta=1$、$\mu=0.9$时，移动平均把教师更新到$0.1$，不是立刻复制成$1$。

无负样本方法避免了显式的负样本队列或大规模两两比较，却仍依赖增强、架构和优化。对比、蒸馏与重建可以组合，但新增目标会引入权重与梯度冲突。应先验证单一目标，再判断组合是否保留了不同的有效信息，而不是因目标数量增加就推断表示更通用。

\section{数据组织与领域继续预训练}
预训练数据应登记来源、设备、时间、工况和使用许可。没有标签不等于可以随意使用：历史记录可能含有机密信息，也可能已经包含后续测试期的数据。对连续波形，应先按设备或时间段划分，再切窗口；否则同一段波形的重叠片段会落入不同集合，让表示看起来比实际更好。归一化统计和滤波参数也只能用规则允许的数据估计。

领域继续预训练在已有参数上使用领域数据延续学习。领域自适应预训练与任务相关语料预训练在语言任务中已有系统研究\cite{gururangan2020dontstop}。对工业波形，相应思路是在通用信号表示上使用目标设备族的历史无标签数据，但应区分部署前可用历史与未来测试期数据。使用无标签测试输入属于不同的信息使用规则，不能悄悄并入普通归纳评估。

只使用狭窄领域数据继续训练可能遗忘通用能力，也可能过拟合设备身份。可比较较小学习率、较短训练、混合原始数据或限制可训练参数等策略。若采用数据混合，可令来源$s$以概率$q_s$被采样，$\sum_sq_s=1$。按样本量取$q_s$使大来源占主导，均衡来源则会重复小数据集；两者改变了优化分布，需要记录而不能仅报告总样本数。

例如两个来源分别有$900$和$100$个窗口，按样本量抽取对应$q_1=0.9,q_2=0.1$；均衡来源则取$q_1=q_2=0.5$。若总共抽取$1000$次，后一种策略平均各抽$500$次，小来源的单个窗口会被重复使用得更多，但并未产生$500$条独立信息。采样比例需要围绕目标设备群确定，不能只以损失更低为由过度重复容易的来源。

固定算力预算下，模型规模、有效样本量、输入长度和训练步数存在取舍。遮蔽率改变任务难度和部分架构的计算量，队列长度改变负样本覆盖，继续训练改变历史知识的保留。比较时应记录实际处理样本或词元数、去重比例、训练时间与峰值显存。跨领域缩放规律不能直接当作当前工业数据的性能承诺。

预训练产物应同时保存模型、输入规范、数据来源和训练截止时间。投影头是否保留、归一化层统计是否更新、通道顺序如何定义，都会影响复用结果。所谓冻结编码器还应说明是否固定归一化运行统计，否则无参数梯度更新也可能改变模型行为。预训练模型版本应与适用设备和数据使用规则绑定。

\section{工业振动信号的完整学习任务设置}
考虑利用多个设备的历史振动记录预训练编码器，再以有限故障标签进行诊断。首先限定预测对象：窗口状态分类、未来故障预测和故障严重度回归需要保留的信息不同。本例以窗口状态分类为教学任务，输入为窗口内多通道波形及允许获得的运行条件，不假定已有实验结果。

按设备划分预训练与下游评估范围，再按连续记录或运行批次组织窗口。跨设备泛化试验中，测试设备的无标签记录也不进入预训练；若研究目标设备无标签适应，则另设一组明确允许该信息的规则。故障状态标签只能来自可核查记录，不能将未标注窗口全部视为正常。

掩码方案可将波形分成连续片段，只向编码器显示可见片段，再预测被遮蔽片段。对比方案可使用不改变窗口状态的裁剪和经过幅值校准的小噪声，增强幅度在训练与验证范围内选择。若裁剪可能移除唯一故障冲击，两视图便不再可靠共享状态，应限制裁剪或改用其他目标。转速变化明显时，可以在明确的工况条件下比较，避免把转速差异一律当作应消除的干扰。

\begin{algorithm}[htbp]
\caption{工业信号预训练与下游表示评价}
\label{alg:industrial-pretraining}
\begin{algorithmic}[1]
\Require 按设备和时间划分的数据、预训练目标、有限下游标签
\State 固定允许使用的数据、窗口边界、归一化与增强规则
\State 初始化编码器及预训练辅助头
\For{每个预训练批次}
 \State 按设备和工况采样窗口，生成掩码或增强视图
 \State 计算选定的重建、对比或自蒸馏损失
 \State 更新在线参数，并按算法更新教师或队列
 \State 监测目标损失、表示方差和不同来源的训练贡献
\EndFor
\State 在固定标签子集上训练线性探测头，并单独进行微调对照
\State 用验证集选择配置，冻结选择后在独立测试集上评价
\end{algorithmic}
\end{algorithm}

对照应包括随机初始化、适用的监督预训练，以及在相同数据范围下的代表性自监督目标。线性探测与全量微调分别报告，至少在多个标签预算下绘制学习曲线，比较每类故障召回、宏平均指标、误报及校准。若使用同一个预训练模型服务多个任务，可另外报告摊销成本，但不能把预训练计算完全隐藏。

消融应保持数据和计算条件尽量一致。例如固定编码器比较随机点遮蔽与连续片段遮蔽，检验收益是否来自更长上下文；固定增强比较是否使用投影头，检验下游表示与训练目标空间的区别；固定预训练样本量比较来源混合，检验模型是否依赖某个设备族。注意单纯匹配训练轮数不一定匹配算力，尤其是不同遮蔽编码和对比批次设置之间。

预训练损失下降但线性探测不改善，可能说明任务主要学到了插值或身份线索。线性探测改善而微调收益消失，则需检查标签量、适应策略和优化预算。某些设备收益明显、另一些退化时，应报告子群差异，而非只用平均分掩盖负迁移。表示可视化只提供诊断线索，不能代替独立分类或预测评价。

监督预训练、自监督目标和表示坍缩可结合下面的教程复习，分别关注训练目标与对比/非对比表示学习。
\section*{辅助学习材料}
\begingroup
\emergencystretch=3em
\begin{itemize}
\item \textbf{Self-Supervised Representation Learning}；作者/平台：Lilian Weng；英文；技术教程。重点阅读 pretext tasks、上下文预测和 Contrastive Predictive Coding，帮助理解本章的重建、对比目标和表示质量。\href{https://lilianweng.github.io/posts/2019-11-10-self-supervised/}{在线阅读}。
\item \textbf{Transfer Learning for Computer Vision}；作者/平台：PyTorch 官方教程；英文；代码教程。重点看 pretrained weights、fine-tuning、feature extraction 和 augmentation，帮助理解本章从通用预训练到领域继续训练的实验边界。\href{https://docs.pytorch.org/tutorials/beginner/transfer\_learning\_tutorial.html}{在线阅读}。
\item \textbf{通俗易懂讲AI--自监督学习}；知乎；中文解说。区分预训练辅助目标与下游目标，比较遮蔽、重建和对比学习；标题与主题据必应索引，原页受限。\href{https://zhuanlan.zhihu.com/p/676130861}{在线阅读}。
\item \textbf{【深度学习】自监督学习详解（self-supervised learning）}；CSDN；中文博客。阅读对比学习和生成学习的入门说明，表示质量仍须由独立下游任务评价；已读取公开页面引言与方法概述。\href{https://blog.csdn.net/qq\_51392112/article/details/128806998}{在线阅读}。
\item \textbf{2021 - 自监督式学习 (二) – BERT简介}；李宏毅；bilibili；中文课程。选看第72集，以语言遮蔽预训练为例对照本章自监督目标。2026年10月4日官方公开元数据显示整套课程累计播放约299.5万，并非本分集播放量；已核对目录与计数，未逐帧观看。\href{https://www.bilibili.com/video/BV1Wv411h7kN/?p=72}{观看分集}。
\end{itemize}
\endgroup
\paragraph{本章小结。}
预训练的设计由三件事共同决定：哪些数据可以使用、什么变化应当保持一致、什么信息必须保留。监督目标利用已有标签，掩码与生成目标恢复缺失或后续内容，对比目标定义样本关系，自蒸馏和统计约束抑制表示退化。这些目标的数学结构对应不同偏好：平方误差输出条件均值，对比交叉熵放大难区分候选的作用，停止梯度改变优化路径，方差与协方差项分别约束常量坍缩和维度重复。理解这些偏好，比记住某个模型名称更有助于判断故障信息会不会被保留。

判断预训练是否有用，要在独立任务上给出相同标签预算，再比较线性探测或微调后的结果。报告中应写清用了哪些历史数据，怎样增强或遮蔽，表示有没有退化，以及故障识别究竟改善多少。辅助损失下降只能说明辅助任务学得更好。下一章接着检查这些表示换到新设备后是否仍适用，并讨论重加权、表示对齐、参数适配和负迁移。
